\clearpage
\fancyhead[L]{\small CERES. ELEMENTS XII.}
\hypertarget{native-304-lower}{}
\begin{center}CERES. ELEMENTS XII.
\end{center}

\begin{center}
\fontsize{10}{12}\selectfont\begin{tabular}{c}
\hline
Göttingische gelehrte Anzeigen. Issue 40. Pp. 393..395. 10 March 1808. \\
\hline
\end{tabular}
\end{center}

The results of last year’s observations of the new planets, Vesta, Pallas and Juno, are already in the hands of the astronomical community. Those concerning Juno were published in these pages only recently. Ceres alone remains, and Prof. \textsc{Gauss} has just finished his investigations of it. In addition to the observations made by Prof. \textsc{Harding} at our observatory in April and May 1807 and some observations from Berlin, Prof. \textsc{Gauss} was able to use a fine series of Milan observations communicated in manuscript. All these observations gave the following result for the latest opposition:


\begin{center}\fontsize{10}{12}\selectfont\setlength{\tabcolsep}{3pt}\renewcommand{\arraystretch}{1.15}\begin{tabular}{lr}
Time of opposition 1807, May 3. .\;.\; & $4^{\text{u}}\,12^{\text{m}}\,44^{\text{s}}$\\
\multicolumn{2}{r}{Göttingen meridian}\\
True longitude .\;.\;.\;.\;.\;.\;.\;,\;.\; & $222^\circ\;14'\;9''1$\\
True geocentric latitude .\;.\;.\;.\; & $10\;40\;15.9$ north.\\
\end{tabular}\end{center}

\noindent This is the \emph{fifth} opposition observed since the discovery of \emph{Ceres}.

\par\indent The previous elements, also published in these pages (1806, p. 1946) [by \textsc{Harding}], differ from this result by only 3 minutes in longitude and only 13 seconds in heliocentric latitude. Small though it is, the discrepancy decisively proves that the perturbations to which Prof. \textsc{Gauss} has so far confined himself no longer suffice to represent all the available observations accurately. So far only perturbations due to Jupiter have been considered, and even among these, those depending on higher powers of the eccentricities have been left out. To reconcile past and future observations accurately, these latter perturbations in particular will therefore have to be taken into account, though for more than one reason it will be advisable to wait for one or a couple of new oppositions before doing this work. For the time being, then, Prof. \textsc{Gauss} has been content to improve the previous elements using only the perturbations in longitude and in the radius vector tabulated in the March 1803 issue of the Monatliche Correspondenz, but with more fully calculated perturbations in latitude (which have yet to be printed), so as to distribute the differences as far as possible across all five oppositions. The following new elements (No. XII in the series) are the result of this investigation. As will be seen, they differ very little from the previous ones. The epochs are for the Göttingen meridian.

\Needspace{18\baselineskip}
\SourceLeaf{305}
\begin{center}\fontsize{10}{12}\selectfont\setlength{\tabcolsep}{3pt}\renewcommand{\arraystretch}{1.15}\begin{tabular}{lrr}
& Mean longitude & Perihelion\\\hline
1802 & $155^\circ\;28'\;15''7$ & $146^\circ\;30'\;22''1$\\
1803 & $233\;38\;7.1$ & $32\;23.3$\\
1804 & $312\;0\;49.5$ & $34\;24.9$\\
1805 & $30\;10\;40.9$ & $36\;26.1$\\
1806 & $108\;20\;32.3$ & $38\;27.4$\\
1807 & $186\;30\;23.7$ & $40\;28.6$\\
1808 & $264\;53\;6.1$ & $42\;30.2$\\
\end{tabular}\end{center}

\begin{center}\fontsize{10}{12}\selectfont\setlength{\tabcolsep}{3pt}\renewcommand{\arraystretch}{1.15}\begin{tabular}{lr}
Daily mean motion .\;.\;.\;.\;.\; & $770''9354$\\
Eccentricity 1806 .\;.\;.\;.\;.\;.\;.\;.\;.\; & $0.0785251$\\
Annual decrease .\;.\;.\;.\;.\;.\;.\;.\; & $0.00000583$\\
Logarithm of the semimajor axis .\;.\;.\; & $0.4420439$\\
Ascending node 1806 .\;.\;.\;.\; & $80^\circ\;53'\;24''6$\\
Annual motion .\;.\;.\;.\;.\;.\;.\;.\; & $+1''48$\\
Inclination of the orbit 1806 .\;.\;.\;.\;.\;.\; & $10^\circ\;37'\;33''2$\\
Annual decrease .\;.\;.\;.\;.\;.\;.\;.\; & $-0''44$\\
\end{tabular}\end{center}

\par\indent These elements agree with the five oppositions observed so far as follows:

\begin{center}\fontsize{9}{11}\selectfont\setlength{\tabcolsep}{3pt}\renewcommand{\arraystretch}{1.15}\begin{tabular}{lrrrr}
& \multicolumn{2}{c}{Difference} & \multicolumn{2}{c}{\shortstack{From\\Elements XIII}}\\
& \shortstack{Heliocentric longitude\\in the orbit} & \shortstack{Heliocentric\\latitude} & & \\\hline
1802 & $-24''1$ & $+1''2$ & $-0''5$ & $-1''9$\\
1803 & $+40.9$ & $-9.6$ & $+22.1$ & $-5.3$\\
1804 & $-21.0$ & $-3.1$ & $-34.9$ & $-1.2$\\
1806 & $+7.8$ & $0$ & $+26.7$ & $-3.6$\\
1807 & $-9.5$ & $-11.3$ & $-26.9$ & $-11.7$\\
1808 & $+62.7$ & & $+20.8$ & $-6.0$\\
\end{tabular}\end{center}

\noindent [\textsc{Gauss} added by hand the comparison with the opposition of 1808 and the differences for Elements No. XIII. \textsc{Schering}.]

